\usepackage{amsmath,amssymb,bm}
\usepackage{booktabs}
\usepackage{graphicx}
\usepackage{hyperref}
\usepackage{cleveref}
\usepackage{xcolor}
\usepackage{enumitem}
\usepackage{microtype}
\usepackage{geometry}
\usepackage{caption}
\usepackage{subcaption}
\usepackage[backend=biber,style=ieee,sorting=none]{biblatex}

\geometry{a4paper,margin=25mm}
\graphicspath{{../system_overview/assets/final/}{../system_overview/assets/selected/}}
\addbibresource{references.bib}
\hypersetup{colorlinks=true,linkcolor=blue!55!black,citecolor=blue!55!black,urlcolor=blue!55!black}

\newcommand{\systemname}{\textsc{PRE-MAP-VLN}}
\newcommand{\pose}{\bm{q}}
\newcommand{\todo}[1]{\textcolor{red!70!black}{[TODO: #1]}}

\newcommand{\overviewfigure}{%
  \begin{figure*}[t]
    \centering
    \IfFileExists{../system_overview/assets/final/system_overview.pdf}{%
      \includegraphics[width=\textwidth]{system_overview.pdf}%
    }{%
      \fbox{\parbox[c][45mm][c]{0.94\textwidth}{\centering
      System Overview placeholder\\
      Export the final figure to\\
      \texttt{paper/system\_overview/assets/final/system\_overview.pdf}}}%
    }
    \caption{Overview of \systemname. The system first constructs a reusable 3D semantic map, then jointly reasons about language constraints, executable observation poses, and collision-free flight paths. Online observations trigger conditional replanning or semantic recovery when needed.}
    \label{fig:overview}
  \end{figure*}%
}
